\section{The canonical adelic line in the remaining slope problem}
\label{sec:AR10}

\begin{theorem}[AR-10: an explicit adelic representative of the complementary height]
Let $A/K$ have dimension $g=2m-1$, let $L$ be a symmetric ample line
bundle rigidified at the identity, and equip $L$ with its canonical
adelic metric. Put $d=m-1$ and
\[
 \overline Q=\mu^*\overline L-p_1^*\overline L-p_2^*\overline L
 \quad\text{on }A\times A,
\]
where $\mu$ is addition. If $z,w\in\operatorname{CH}^m_{(1)}(A)$
are primitive in the Chow sense $Lz=Lw=0$, then
\begin{equation}\label{eq:AR10}
 (-1)^m h_K(z,w)=
 \frac{(-1)^{m-1}}{(2m-1)!\,h^0(A,L)\,[K:\mathbf Q]}
 \widehat{\deg}\bigl(\overline Q^{\,2m-1}\cdot(z\times w)\bigr).
\end{equation}
Here $h_K$ is the normalized raw Beilinson--Bloch pairing of the
arithmetic section. The arithmetic intersection on the right uses
rational cycle representatives and canonical metrics.

For a chosen rational representative $z=\sum_i a_i[V_i]$, define the
integrable adelic rational line on $\operatorname{Spec}K$
\[
 \overline{\mathcal B}_z=
 \frac{(-1)^{m-1}}{(2m-1)!h^0(A,L)}
 \sum_{i,j}a_i a_j
 \bigl\langle\underbrace{\overline Q,\ldots,\overline Q}_{2m-1}
 \bigr\rangle_{(V_i\times V_j)/K}.
\]
Its normalized arithmetic degree is $(-1)^m h_K(z,z)$. Thus for
$z\in\operatorname{CH}^m(N)$ in AR-2, the line required as the source
of the slope criterion AR-6 has an explicit construction from the
cycle and the canonical polarization metric. The residual requirement
of that criterion is an injective map from its dual to a hermitian
bundle satisfying the stated maximal-slope and operator-norm bound.
\end{theorem}

\begin{proof}
We use Gu, \emph{The Beilinson--Bloch heights of generalized Ceresa
cycles}, arXiv:2609.33223v2 (2 October 2026), Theorem 1.3, printed p.~2,
and Proposition 5.3, printed pp.~12--13. These are cited formulas;
the specialization and the slope-line construction are the deduction
here. Gu uses the unnormalized height before division by $[K:\mathbf Q]$,
so we keep this division explicitly on arithmetic intersections.

A grade-one codimension-$m$ class has eigenvalue $(-1)^{2m-1}=-1$
under $[-1]^*$. Since $[-1]$ is an involution, pushforward has the same
action. Hence the normalized generalized Ceresa cycle
$\operatorname{Ce}(z)=(z-[-1]_*z)/2$ is exactly $z$.
Its pure dimension is $d=g-m=m-1$. Gu's Theorem 1.3 therefore gives
\[
 h_{\mathrm{BB},F}(z)=
 \frac{(-1)^d}{(2d+1)!\,[K:\mathbf Q]}
 \widehat{\deg}\bigl(\overline Q^{2d+1}\cdot(z\times z)\bigr).
\]
For the primitive grade-one class $z$, Proposition 5.3 gives
\[
 h_{\mathrm{BB},F}(z)=
 \frac{h^0(A,L)}{(g+1-2m)!}\,h_{\mathrm{BB},\mathrm{Lef}}(z)
 =h^0(A,L)(-1)^m h_K(z,z).
\]
Indeed $g+1-2m=0$, and the Lefschetz operator in that height is the
identity. Combining the two identities proves the diagonal case of
\eqref{eq:AR10}. The kernel $\overline Q$ is invariant under interchange
of the factors, so its arithmetic intersection against $z\times w$
is symmetric in $z,w$. Apply the diagonal identity to $z+w$ and
subtract the identities for $z$ and $w$; division by two proves the
bilinear statement.

The relative dimension of each $V_i\times V_j$ is $2d=2m-2$.
The displayed Deligne pairing therefore has exactly one more line-bundle
argument than the relative dimension. Its arithmetic degree is the
arithmetic intersection of those arguments, including their canonical
metrics. This is the metric Deligne-pairing identity recalled in Gu,
\S2.2, printed p.~5; its extension to rational cycles is multilinearity,
as also stated in Corollary 1.6, printed p.~3. It proves the assertion
about $\overline{\mathcal B}_z$. Finite base change multiplies the degree
by the extension degree, so the normalized equality is unchanged.
The arithmetic degree is independent of the chosen representative by
\eqref{eq:AR10}; no canonical isometry between the lines from all
representatives is needed or asserted.

Finally, use the following adelic rank-one version of the slope
argument AR-6. An integrable adelic rational line is not automatically
an integral hermitian lattice, so this extension is made explicit.
Let $\overline E=\overline{\mathcal B}_z^{\vee}$, let
$\varphi:E\hookrightarrow F$ be an injective linear map of adelic
normed spaces, and give its image $I$ the norms induced from $F$.
With normalized local weights $n_v$, put
$h(\varphi)=\sum_v n_v\log\|\varphi\|_v$.
For a nonzero vector $e\in E$, the local operator bound yields
$\|\varphi(e)\|_v\le\|\varphi\|_v\|e\|_v$.
Negate logarithms and sum to get
\[
 \widehat{\deg}_n E\le
 \widehat{\deg}_n I+h(\varphi)
 \le\widehat\mu_{\max}(F)+h(\varphi).
\]
The product formula makes the degrees independent of the chosen vector.
This proof works for integrable adelic norms and rational powers by
linearity; no passage to an integral lattice is being assumed.
Since $\widehat{\deg}_n E=-(-1)^m h_K(z,z)$, a nonpositive
right side proves the required nonnegative height, and a negative
right side proves strict positivity. The construction has fixed its metric through
$\overline L$; it has not assigned an arbitrary metric using the
unknown height. Therefore the only additional hypotheses of that
argument are the actual injective map and its norm/slope estimate.
\end{proof}

\paragraph{Why this does not already give the sign.}
The adelic kernel $\overline Q$ is integrable, not nef in general.
Even its underlying line bundle has
\[
 \Delta^*Q=2L,\qquad (1,-1)^*Q=-2L
 \quad\text{in }\operatorname{NS}(A)_{\mathbf Q}.
\]
Both identities follow from the theorem of the square and the
rigidification at the identity. The second restricts negatively to
curves. Consequently one cannot replace the remaining arithmetic
inequality by a blanket assertion that intersections of $\overline Q$
are nonnegative. Gu's own Conjecture 1.1, printed p.~1, explicitly leaves
the general Fourier-height positivity problem open. The construction
above supplies the geometric height line for a Bost approach; it does
not supply the norm estimate.
